\usepackage[utf8]{inputenc}
\usepackage[T1]{fontenc}
\usepackage{textcomp}
\usepackage{scalerel}
\usepackage{marginnote}
\usepackage[light]{inter}
\usepackage{url}
\usepackage[unicode]{hyperref}
\hypersetup{
    colorlinks=true,
    linkcolor=blue,
    filecolor=magenta,      
    urlcolor=cyan,
    pdftitle={ST202 Notes},
    }

\usepackage{amsmath, amsfonts, mathtools, amsthm, amssymb}
\usepackage{mathrsfs}
\usepackage{cancel}
\usepackage{enumerate}
\usepackage{authblk}

\usepackage{graphicx}
\usepackage{float}
\usepackage[usenames, dvipsnames]{xcolor}

\def\🎉{\scalerel*{\includegraphics{figures/1f389.eps}}{\textrm{\textbigcircle}}}
\let\implies\Rightarrow
\let\impliedby\Leftarrow
\let\iff\Leftrightarrow
\let\epsilon\varepsilon

\newcommand\hr{
    \noindent\rule[0.5ex]{\linewidth}{0.5pt}
}

\usepackage{tikz}
\usepackage{tikz-cd}

% theorems
\usepackage{thmtools}
\usepackage[framemethod=TikZ]{mdframed}
\mdfsetup{skipabove=1em,skipbelow=0em, innertopmargin=5pt, innerbottommargin=6pt}

\theoremstyle{definition}

\makeatletter

\declaretheoremstyle[headfont=\bfseries\rmfamily, bodyfont=\normalfont, mdframed={ nobreak } ]{thmgreenbox}
\declaretheoremstyle[headfont=\bfseries\rmfamily, bodyfont=\normalfont, mdframed={ nobreak } ]{thmredbox}
\declaretheoremstyle[headfont=\bfseries\rmfamily, bodyfont=\normalfont]{thmbluebox}
\declaretheoremstyle[headfont=\bfseries\rmfamily, bodyfont=\normalfont]{thmblueline}
\declaretheoremstyle[headfont=\bfseries\rmfamily, bodyfont=\normalfont, numbered=no, mdframed={ rightline=false, topline=false, bottomline=false, }, qed=\qedsymbol ]{thmproofbox}
\declaretheoremstyle[headfont=\bfseries\rmfamily, bodyfont=\normalfont, numbered=no, mdframed={ nobreak, rightline=false, topline=false, bottomline=false } ]{thmexplanationbox}

\declaretheoremstyle[headfont=\bfseries\rmfamily, bodyfont=\normalfont, numbered=no, mdframed={ nobreak, rightline=false, topline=false, bottomline=false } ]{thmexplanationbox}


\declaretheorem[numberwithin=section, style=thmgreenbox, name=Definition]{definition}
\declaretheorem[sibling=definition, style=thmredbox, name=Corollary]{corollary}
\declaretheorem[sibling=definition, style=thmredbox, name=Proposition]{prop}
\declaretheorem[sibling=definition, style=thmredbox, name=Theorem]{theorem}
\declaretheorem[sibling=definition, style=thmredbox, name=Lemma]{lemma}



\declaretheorem[sibling=definition, style=thmbluebox,  name=Example]{eg}
\declaretheorem[sibling=definition, style=thmbluebox,  name=Counterexample]{cx}
\declaretheorem[sibling=definition, style=thmblueline, name=Remark]{remark}
\declaretheorem[numbered=no, style=thmexplanationbox, name=Proof]{explanation}
\declaretheorem[numbered=no, style=thmproofbox, name=Proof]{replacementproof}
\declaretheorem[style=thmbluebox, numbered=no, name=Exercise]{ex}
\declaretheorem[style=thmblueline, numbered=no, name=Note]{note}

\renewenvironment{proof}[1][\proofname]{\begin{replacementproof}}{\end{replacementproof}}

\AtEndEnvironment{eg}{\null\hfill$\diamond$}%

\newtheorem*{notation}{Notation}
\newtheorem*{recall}{Recall}
\newtheorem*{problem}{Problem}
\newtheorem*{observe}{Observe}
\newtheorem*{property}{Property}
\newtheorem*{intuition}{Intuition}

\usepackage{etoolbox}

 
\usepackage{xifthen}

\def\thm@space@setup{
\thm@preskip=\parskip \thm@postskip=0pt}

\def\testdateparts#1{\dateparts#1\relax}
\def\dateparts#1 #2 #3 #4 #5\relax{
    \marginnote{\hsize3cm\small\textsf{\mbox{#1 #2 #3 #5}}}
}

\def\@lecture{}%
\newcommand{\lecture}[4]{
    \ifthenelse{\isempty{#4}}{%
        \def\@lecture{Week #1: #2}%
    }{%
        \def\@lecture{Week #1: #2 $\,\mid\,$  #4}%
    }%
    \section{\@lecture}
	\testdateparts{#3}
}

\usepackage{nameref}
\usepackage{fancyhdr}
\pagestyle{fancy}

\renewcommand{\chaptermark}[1]{ \markboth{#1}{} }

\fancyhead[L]{\small\textsf\@lecture}
\fancyhead[R]{\small\textsf{\thechapter} \normalsize $ \ \diagup \ $ \small\textsf{\nouppercase\leftmark}}
\fancyfoot[C]{\thepage}
\renewcommand{\headrulewidth}{0pt}

\makeatother

%Figure Support
\usepackage{import}
\pdfminorversion=7
\usepackage{pdfpages}
\usepackage{transparent}

\newcommand{\incfig}[3][0.8]{%
	\begin{figure}[H]
	\centering
    \def\svgwidth{#1\columnwidth}
    \import{./figures/}{#2.pdf_tex}
	\ifthenelse{\isempty{#3}}{%
	}{%
		\caption{#3}%
	}%
	\label{fig:#2}
	\end{figure}
}

\pdfsuppresswarningpagegroup=1

\linespread{1.2}

\author{Tay Meshkinyar \vskip.2cm Dr. Milt Mavrakakis}

\DeclareRobustCommand\iff{\;\Longleftrightarrow\;}
\newcommand\independent{\protect\mathpalette{\protect\independenT}{\perp}}
\def\independenT#1#2{\mathrel{\rlap{$#1#2$}\mkern2mu{#1#2}}}

\newcommand{\ind}{\independent}
\newcommand{\RR}{\mathbb R}
\newcommand{\FF}{\mathbb F}
\newcommand{\LL}{\mathcal L}
\newcommand{\lL}{\mathcal l}
\newcommand{\NN}{\mathbb N}
\newcommand{\EE}{\mathbb E}
\newcommand{\hh}{\mathbf h}
\newcommand{\HH}{\mathbf H}
\newcommand{\bS}{\mathbf S}
\newcommand{\bs}{\mathbf s}
\newcommand{\bT}{\mathbf T}
\newcommand{\bt}{\mathbf t}
\newcommand{\bee}{\mathbf b}
\newcommand{\ba}{\mathbf a}
\newcommand{\ee}{\mathbf e}
\newcommand{\xx}{\mathbf x}
\newcommand{\XX}{\mathbf X}
\newcommand{\yy}{\mathbf y}
\newcommand{\YY}{\mathbf Y}
\newcommand{\uu}{\mathbf u}
\newcommand{\UU}{\mathbf U}
\newcommand{\bv}{\mathbf v}
\newcommand{\gee}{\mathbf g}
\newcommand{\intR}{\int_{\mathbb R}}
\newcommand{\intRR}{\int_{\mathbb R^2}}
\newcommand{\suminf}[1]{\sum^{\infty}_{#1=0}}
\newcommand{\sumN}[1]{\sum_{#1\in \mathbb N}}
\newcommand{\sumNZ}[1]{\sum_{#1\in \mathbb N_0}}
\newcommand{\pder}[2][]{\frac{\partial#1}{\partial#2}}
\newcommand{\der}[2][]{\frac{d#1}{d#2}}
\newcommand{\Det}[4]{\left| \
\begin{matrix}
#1 & #2 \\[1ex]
#3 & #4  
\end{matrix}\
\right|}

\newcommand{\fXY}{f_{X,Y}(x,y)}
\newcommand{\fYy}{f_{\mathbf Y}(\mathbf y)}
\newcommand{\fUV}{f_{U,V}(u,v)}
\newcommand{\FXY}{F_{X,Y}(x,y)}
\newcommand{\FUV}{F_{U,V}(u,v)}
\newcommand{\MXY}{M_{X,Y}(t,u)}
\newcommand{\KXY}{K_{X,Y}(t,u)}
\newcommand{\fYgX}{f_{Y\mid X}(y\mid x)}
\newcommand{\FYgX}{F_{Y\mid X}(y\mid x)}
\newcommand{\MYX}{M_{Y\mid X}(u\mid X)}
\newcommand{\MYx}{M_{Y\mid X}(u\mid x)}
\DeclareMathOperator*{\argmax}{arg\,max}
\DeclareMathOperator*{\argmin}{arg\,min}
\DeclareMathOperator*{\Bias}{Bias}
\DeclareMathOperator*{\MSE}{MSE}
\DeclareMathOperator*{\Var}{Var}
\DeclareMathOperator*{\Cov}{Cov}
\DeclareMathOperator*{\Corr}{Corr}
\DeclareMathOperator*{\Exp}{Exp}
\DeclareMathOperator*{\Normal}{Normal}
\DeclareMathOperator*{\Bin}{Bin}
\DeclareMathOperator*{\Bernoulli}{Bernoulli}
\DeclareMathOperator*{\Poisson}{Poisson}
\DeclareMathOperator*{\Unif}{Unif}
\DeclareMathOperator*{\Beta}{Beta}
\DeclareMathOperator*{\Geo}{Geo}
\newcommand{\YOne}{{Y_{(1)}}}
\newcommand{\Yn}{{Y_{(n)}}}
\newcommand{\Yi}{{Y_{(i)}}}
\newcommand{\Yii}[1][]{{Y_{(i#1)}}}
\newcommand{\btheta}{{\boldsymbol{\theta}}}
\newcommand{\bbeta}{{\boldsymbol{\beta}}}
\newcommand{\bepsilon}{{\boldsymbol{\epsilon}}}
\newcommand{\DD}{\mathcal{D}}

